\chapter{题型/关键词 → 算法索引}\label{ux9898ux578bux5173ux952eux8bcd-ux7b97ux6cd5ux7d22ux5f15}

本页专门给``现场查名字''使用。下表中的算法先写完整中文/英文名称，括号内才是题解和洛谷题面中常见的简称；不要只记简称。标注``作用''的条目为 2100 分及以上的重点内容。

\section{一、博弈论}\label{ux4e00ux535aux5f08ux8bba}

\begin{longtable}[]{@{}
  >{\raggedright\arraybackslash}p{(\linewidth - 4\tabcolsep) * \real{0.3333}}
  >{\raggedright\arraybackslash}p{(\linewidth - 4\tabcolsep) * \real{0.3333}}
  >{\raggedright\arraybackslash}p{(\linewidth - 4\tabcolsep) * \real{0.3333}}@{}}
\toprule\noalign{}
\begin{minipage}[b]{\linewidth}\raggedright
约分
\end{minipage} & \begin{minipage}[b]{\linewidth}\raggedright
完整名称（常见写法）
\end{minipage} & \begin{minipage}[b]{\linewidth}\raggedright
2100 分及以上：作用
\end{minipage} \\
\midrule\noalign{}
\endhead
\bottomrule\noalign{}
\endlastfoot
1800--2000 & 必胜态/必败态动态规划（Win/Lose Dynamic Programming） & --- \\
1800--2100 & 有向无环图上的 Sprague--Grundy 函数与最小未出现值（Directed Acyclic Graph、SG、mex） & 把无偏无环博弈状态转成整数，便于判断胜负或合并子游戏。 \\
1800--2100 & 多个独立子博弈的 Sprague--Grundy 值异或 & 判断游戏能否拆成互不影响的部分，并用异或合并各部分。 \\
1800--2000 & Bash 取石子游戏及限制取石子变形 & --- \\
1800--2000 & 普通 Nim 游戏的必胜操作构造 & --- \\
1900--2100 & 反常 Nim 游戏（Misère Nim） & --- \\
1900--2200 & 阶梯 Nim 游戏（Staircase Nim） & 只对特定编号台阶取异或，处理石子只能向下移动的游戏。 \\
1900--2200 & 图游戏建模（Graph Game Modeling） & 把棋子移动、可达性和取点过程转成状态图，再做胜负或 Grundy 分析。 \\
2000--2300 & Sprague--Grundy 函数打表与周期发现 & 将超大取石子数量压缩到前导段和周期内，避免按数量线性计算。 \\
2100--2400 & 大范围取石子游戏的 Sprague--Grundy 周期压缩 & 对超大状态值利用已证明的周期直接回答，常用于 \passthrough{\lstinline!n!} 达到 \passthrough{\lstinline!10\^18!} 的题。 \\
1900--2200 & 普通 Nim 游戏的必胜策略恢复 & 根据异或和构造具体取哪一堆、取多少个。 \\
2000--2300 & 状态压缩博弈动态规划（State-Compressed Game Dynamic Programming） & 用位掩码保存少量棋子/资源状态，枚举合法操作并记忆化。 \\
1800--2100 & 记忆化博弈搜索（Memoized Game Search） & 对有大量重复状态的博弈搜索去重，避免重复递归。 \\
1900--2200 & 极大极小搜索（Minimax Search） & 双方操作不同或不能拆成独立子游戏时，在完整博弈树上求最优决策。 \\
2000--2300 & Alpha-Beta 剪枝（Alpha-Beta Pruning） & 在极大极小搜索中利用上下界跳过不可能改变答案的分支。 \\
2100--2400 & 博弈树状态去重与状态哈希（Transposition Table and State Hashing） & 合并不同走法到达的同一局面，降低棋类搜索的状态数。 \\
2200--2500 & 多阶段/带资源限制的组合游戏 & 处理多个阶段、共享资源和操作规则变化，不能直接套独立游戏异或。 \\
2300--2500 & 有环图上的博弈状态分析（Cyclic Graph Game Analysis） & 区分必胜、必败和平局，处理会回到旧状态的图博弈。 \\
\end{longtable}

查阅位置：博弈专题基础模板后紧跟的博弈论详解与自测样例。

\section{二、高级数据结构}\label{ux4e8cux9ad8ux7ea7ux6570ux636eux7ed3ux6784}

\begin{longtable}[]{@{}
  >{\raggedright\arraybackslash}p{(\linewidth - 4\tabcolsep) * \real{0.3333}}
  >{\raggedright\arraybackslash}p{(\linewidth - 4\tabcolsep) * \real{0.3333}}
  >{\raggedright\arraybackslash}p{(\linewidth - 4\tabcolsep) * \real{0.3333}}@{}}
\toprule\noalign{}
\begin{minipage}[b]{\linewidth}\raggedright
约分
\end{minipage} & \begin{minipage}[b]{\linewidth}\raggedright
完整名称（常见写法）
\end{minipage} & \begin{minipage}[b]{\linewidth}\raggedright
2100 分及以上：作用
\end{minipage} \\
\midrule\noalign{}
\endhead
\bottomrule\noalign{}
\endlastfoot
1800--2100 & 高阶前缀和与多阶差分（Higher-Order Prefix Sums and Differences） & --- \\
1900--2200 & 多重懒标记线段树（Segment Tree with Multiple Lazy Tags） & 同时维护多种区间修改，并正确合成和下传标记。 \\
1900--2200 & 仿射标记线段树（Affine-Transformation Segment Tree） & 统一处理区间乘法、区间加法和区间和，标记形式为 \passthrough{\lstinline!x → ax+b!}。 \\
1800--2100 & 线段树上的二分（Segment Tree Binary Search） & 在区间最大值、容量或非负前缀和上找第一个满足条件的位置。 \\
1800--2100 & 莫队算法（Mo's Algorithm） & --- \\
2200--2400 & 带修改莫队算法（Mo's Algorithm with Modifications） & 离线处理区间询问和时间维修改，状态移动为左端点、右端点、修改时间。 \\
2200--2400 & 回滚莫队算法（Rollback Mo's Algorithm） & 在不支持普通删除或需要撤销的区间统计中，用日志回退维护状态。 \\
2000--2300 & 三维偏序的 CDQ 分治（Three-Dimensional Dominance CDQ Divide and Conquer） & 把一个维度分治、一个维度排序、最后一维交给树状数组统计。 \\
2200--2400 & CDQ 分治优化动态规划（CDQ Divide and Conquer Optimization） & 让左半状态只贡献给右半状态，优化具有时间/坐标偏序的动态规划转移。 \\
1900--2200 & 可持久化字典树（Persistent Trie / Persistent Prefix Tree） & 对历史版本和区间版本做异或最大值、第 \passthrough{\lstinline!k!} 大等查询。 \\
2200--2400 & 可持久化并查集（Persistent Disjoint Set Union） & 保存每个时刻的连通关系，支持复制历史版本、合并和查询连通性。 \\
1900--2200 & 树上启发式合并（Disjoint Set Union on Tree，DSU on Tree） & 在线性树上统计每个子树的颜色频次、众数或频率函数。 \\
2200--2500 & 长链剖分（Long-Chain Decomposition） & 用最长向下链优化树上深度 DP、\passthrough{\lstinline!k!} 级祖先和同深度信息合并。 \\
2100--2400 & 点分树/重心树（Centroid Tree） & 把树上距离查询转成重心祖先链上的若干距离容器，支持在线修改和查询。 \\
2100--2400 & 李超线段树（Li Chao Segment Tree） & 动态维护直线集合，在指定横坐标查询最大值或最小值。 \\
2200--2400 & 线段插入版李超线段树（Li Chao Segment Tree for Line Segments） & 让每条直线只在一个横坐标区间有效，处理线段插入与在线最值查询。 \\
2300--2500 & 区间最值势能线段树（Segment Tree Beats） & 维护区间取最小/最大与区间和，利用最大值、次大值和势能保证均摊复杂度。 \\
2200--2500 & 伸展树（Splay Tree） & 通过旋转把访问节点移到根，维护动态序列、排名和区间翻转。 \\
2200--2400 & 文艺平衡树/序列伸展树（Reversible Splay Tree） & 用伸展树维护序列区间反转、插入、删除和区间查询。 \\
2400--2500 & Link-Cut Tree 动态树（Link-Cut Tree，LCT） & 维护动态森林的连边、断边、换根和路径信息。 \\
2400--2500 & 动态动态规划与矩阵线段树（Dynamic Dynamic Programming and Matrix Segment Tree） & 把局部状态转移写成矩阵，用线段树合并并支持单点修改后的全局答案。 \\
\end{longtable}

\section{三、数学与组合计数}\label{ux4e09ux6570ux5b66ux4e0eux7ec4ux5408ux8ba1ux6570}

\begin{longtable}[]{@{}
  >{\raggedright\arraybackslash}p{(\linewidth - 4\tabcolsep) * \real{0.3333}}
  >{\raggedright\arraybackslash}p{(\linewidth - 4\tabcolsep) * \real{0.3333}}
  >{\raggedright\arraybackslash}p{(\linewidth - 4\tabcolsep) * \real{0.3333}}@{}}
\toprule\noalign{}
\begin{minipage}[b]{\linewidth}\raggedright
约分
\end{minipage} & \begin{minipage}[b]{\linewidth}\raggedright
完整名称（常见写法）
\end{minipage} & \begin{minipage}[b]{\linewidth}\raggedright
2100 分及以上：作用
\end{minipage} \\
\midrule\noalign{}
\endhead
\bottomrule\noalign{}
\endlastfoot
1800--2200 & 高级容斥与倍数容斥（Advanced Inclusion--Exclusion） & --- \\
1900--2200 & 卢卡斯定理（Lucas Theorem） & --- \\
1900--2200 & 扩展中国剩余定理（Extended Chinese Remainder Theorem，扩展 CRT） & 合并模数不互质的同余方程，并处理解的存在性和最小非负解。 \\
2200--2500 & 模合数的组合数（Binomial Coefficients Modulo a Composite Number） & 在不能直接使用阶乘逆元或 Lucas 定理时，按模数质因数分解计算组合数。 \\
1800--2200 & 矩阵快速幂与线性递推（Matrix Exponentiation and Linear Recurrence） & --- \\
2000--2300 & 高斯消元求秩与解空间（Gaussian Elimination for Rank and Solution Space） & 判断方程组是否有解、自由元数量以及线性约束的解空间维数。 \\
2100--2400 & 模意义高斯消元（Gaussian Elimination over a Finite Field） & 在模质数或可逆环上求线性方程组，常用于概率、计数和线性约束。 \\
2100--2400 & 高斯消元求行列式（Determinant by Gaussian Elimination） & 计算矩阵行列式，作为矩阵树定理等组合计数的核心子程序。 \\
2200--2500 & 莫比乌斯反演（Möbius Inversion） & 把``所有倍数/所有约数''的累计统计反推出``恰好倍数/恰好最大公约数''的统计。 \\
2200--2400 & 原根判定与原根求解（Primitive Root Testing and Construction） & 构造模意义下生成整个乘法群的元素，为离散对数和周期问题提供坐标。 \\
2200--2400 & Baby-step Giant-step 离散对数算法（BSGS） & 在模数较小时求 \passthrough{\lstinline!a\^x=b!} 的离散对数，复杂度约为平方根级别。 \\
2300--2500 & 扩展 Baby-step Giant-step 离散对数算法（Extended BSGS） & 处理底数与模数不互质的离散对数，先剥离公因子再进入 BSGS。 \\
2300--2500 & 快速傅里叶变换与数论变换 & 在合适质数模数下快速计算多项式卷积、计数合并和字符串卷积。 \\
2200--2400 & 普通生成函数基础（Ordinary Generating Functions） & 把计数递推和组合对象转成多项式乘法或系数提取问题。 \\
2300--2500 & 整数分拆生成函数（Generating Functions for Integer Partitions） & 计算整数拆分、硬币组合和无序选取的方案数。 \\
2400--2500 & 集合分拆与排列计数生成函数（Generating Functions for Set Partitions and Permutations） & 对集合结构、排列结构和组件合并进行高阶计数。 \\
2000--2300 & Burnside 引理（Burnside's Lemma） & 把等价染色的数量转成群元素固定点数量的平均值。 \\
2200--2500 & Pólya 计数定理（Pólya Enumeration Theorem） & 按置换的循环结构批量计算旋转、翻转对称下的染色方案。 \\
2100--2400 & Legendre 符号（Legendre Symbol） & 快速判断模奇质数下一个数是否为二次剩余。 \\
2200--2500 & Tonelli--Shanks 二次剩余算法（Tonelli--Shanks Algorithm） & 在奇质数模下求平方根，处理 \passthrough{\lstinline!x²=a mod p!} 的有根和无根情况。 \\
2300--2500 & Cipolla 二次剩余算法（Cipolla's Algorithm） & 另一种模奇质数平方根算法，适合需要稳定实现的二次剩余题。 \\
2000--2300 & Prüfer 序列（Prüfer Sequence） & 把标号树编码成长度为 \passthrough{\lstinline!n-2!} 的序列，并利用度数公式计数。 \\
2200--2400 & Prüfer 序列计数应用 & 由点度数、叶子数量和限制条件反推标号树数量。 \\
2200--2500 & 矩阵树定理（Matrix-Tree Theorem） & 用拉普拉斯矩阵的余子式计算无向图生成树数量。 \\
2300--2500 & Lindström--Gessel--Viennot 引理（LGV 引理） & 用行列式计算不相交路径组的数量，处理格点路径和网络路径计数。 \\
1800--2200 & 鸽巢原理构造（Pigeonhole Principle Constructions） & --- \\
1800--2200 & 线性递推的矩阵构造（Matrix Construction for Linear Recurrences） & --- \\
\end{longtable}

\section{四、字符串算法}\label{ux56dbux5b57ux7b26ux4e32ux7b97ux6cd5}

\begin{longtable}[]{@{}
  >{\raggedright\arraybackslash}p{(\linewidth - 4\tabcolsep) * \real{0.3333}}
  >{\raggedright\arraybackslash}p{(\linewidth - 4\tabcolsep) * \real{0.3333}}
  >{\raggedright\arraybackslash}p{(\linewidth - 4\tabcolsep) * \real{0.3333}}@{}}
\toprule\noalign{}
\begin{minipage}[b]{\linewidth}\raggedright
约分
\end{minipage} & \begin{minipage}[b]{\linewidth}\raggedright
完整名称（常见写法）
\end{minipage} & \begin{minipage}[b]{\linewidth}\raggedright
2100 分及以上：作用
\end{minipage} \\
\midrule\noalign{}
\endhead
\bottomrule\noalign{}
\endlastfoot
1800--2000 & Booth 循环字符串最小表示法（Booth's Algorithm） & --- \\
2000--2300 & 回文自动机（Palindromic Automaton，PAM） & 维护每个位置的最长回文后缀和所有不同回文串。 \\
2100--2300 & 回文自动机出现次数统计 & 沿回文自动机失配指针汇总每个回文串的出现次数。 \\
2300--2500 & 回文自动机与动态规划 & 把回文后缀链转成分割、计数或最优值转移。 \\
2100--2300 & 后缀数组、最长公共前缀与区间最小值查询（Suffix Array、Longest Common Prefix、Range Minimum Query） & 对后缀排序后用最长公共前缀和区间最小值查询回答任意后缀/子串关系。 \\
2200--2400 & 字典序第 \passthrough{\lstinline!k!} 个不同子串 & 用后缀数组或后缀自动机统计每个前缀分支的子串数量并恢复答案。 \\
2200--2400 & 多字符串后缀数组（Generalized Suffix Array） & 把多个字符串放在一起比较后缀，求跨字符串最长公共子串等。 \\
2100--2300 & 后缀自动机的结束位置集合树（Suffix Automaton Endpos Tree） & 按后缀链接树汇总出现次数、包含关系和子串贡献。 \\
2000--2300 & 后缀自动机出现次数与不同子串统计（Suffix Automaton） & 线性构建所有子串状态，统计出现次数、不同子串数量和最大贡献。 \\
2300--2500 & 多字符串后缀自动机（Multiple-String Suffix Automaton） & 在多个文本上合并统计每个状态的覆盖串数或共同子串。 \\
2200--2400 & 后缀自动机求字典序第 \passthrough{\lstinline!k!} 个子串 & 在后缀自动机状态有向无环图上做计数并按字符恢复路径。 \\
2200--2400 & Aho--Corasick 多模式匹配自动机的失配树（Aho--Corasick Automaton Fail Tree） & 将文本经过节点的次数沿失配树汇总到模式串，处理后缀包含关系。 \\
2200--2400 & Aho--Corasick 自动机与树状数组 & 用失配树 DFS 序支持模式串动态开关、子树加和查询。 \\
2300--2500 & 动态插入/删除模式串的 Aho--Corasick 自动机 & 将动态模式集合拆成若干静态自动机，支持在线增删和匹配统计。 \\
2300--2500 & 二进制分组动态 Aho--Corasick 自动机 & 用二进制分组控制多个自动机数量，使动态操作达到对数级组数。 \\
1800--2100 & 哈希二分最长公共前缀（Hash-Based Longest Common Prefix） & --- \\
\end{longtable}

\section{五、计算几何}\label{ux4e94ux8ba1ux7b97ux51e0ux4f55}

\begin{longtable}[]{@{}
  >{\raggedright\arraybackslash}p{(\linewidth - 4\tabcolsep) * \real{0.3333}}
  >{\raggedright\arraybackslash}p{(\linewidth - 4\tabcolsep) * \real{0.3333}}
  >{\raggedright\arraybackslash}p{(\linewidth - 4\tabcolsep) * \real{0.3333}}@{}}
\toprule\noalign{}
\begin{minipage}[b]{\linewidth}\raggedright
约分
\end{minipage} & \begin{minipage}[b]{\linewidth}\raggedright
完整名称（常见写法）
\end{minipage} & \begin{minipage}[b]{\linewidth}\raggedright
2100 分及以上：作用
\end{minipage} \\
\midrule\noalign{}
\endhead
\bottomrule\noalign{}
\endlastfoot
1800--2000 & 极角排序（Polar-Angle Sorting） & --- \\
1900--2200 & 射线排序与半平面方向判定（Ray Sorting and Half-Plane Orientation Test） & 用叉积和半平面标记稳定地排序方向、判断左右侧和旋转顺序。 \\
2000--2200 & 旋转卡壳求凸包直径（Rotating Calipers for Convex Polygon Diameter） & 在凸包上在线性时间内求最远点对或最大距离。 \\
2100--2300 & 旋转卡壳求最大三角形（Rotating Calipers for Maximum-Area Triangle） & 在凸多边形上优化三点枚举，求最大面积三角形。 \\
1900--2200 & 凸多边形内点二分（Binary Point Location in a Convex Polygon） & 利用凸性把点内判定从线性扫描降到对数复杂度。 \\
2100--2400 & 凸多边形半平面切割（Convex Polygon Cutting） & 用一条有向直线保留多边形一侧，连续切出新的凸多边形。 \\
2300--2500 & 半平面交（Half-Plane Intersection，HPI） & 求多个线性不等式的公共区域及其面积、顶点和空集情况。 \\
2000--2300 & 最近点对分治（Closest Pair of Points Divide and Conquer） & 在平面点集中求最小欧氏距离，复杂度为 \passthrough{\lstinline!O(n log n)!}。 \\
1900--2200 & 圆与直线交点（Line-Circle Intersection） & --- \\
2000--2200 & 两圆交点（Circle-Circle Intersection） & --- \\
1800--2100 & 点到圆/圆与圆位置关系 & --- \\
2200--2500 & 两圆公切线（Common Tangents of Two Circles） & 求两圆的外公切线、内公切线及切点，处理相切和包含边界。 \\
2100--2300 & 最小圆覆盖（Minimum Enclosing Circle） & 随机增量维护包含所有点的最小圆，期望线性时间完成。 \\
2300--2500 & 多边形与圆面积交（Polygon-Circle Intersection Area） & 将多边形边分解为三角形/圆弓，精确累加相交面积。 \\
2000--2300 & 扫描线求矩形并面积（Sweep Line for Rectangle Union Area） & 通过事件排序和线段树维护当前横向覆盖长度。 \\
2200--2400 & 扫描线求覆盖次数面积（Sweep Line for Area Covered at Least \passthrough{\lstinline!k!} Times） & 在线段树中维护覆盖层数，统计覆盖至少 \passthrough{\lstinline!k!} 次的面积。 \\
1800--2100 & 几何去重与方向向量规范化（Geometric Deduplication and Normalized Direction Vectors） & --- \\
\end{longtable}

\section{六、动态规划进阶}\label{ux516dux52a8ux6001ux89c4ux5212ux8fdbux9636}

\begin{longtable}[]{@{}
  >{\raggedright\arraybackslash}p{(\linewidth - 4\tabcolsep) * \real{0.3333}}
  >{\raggedright\arraybackslash}p{(\linewidth - 4\tabcolsep) * \real{0.3333}}
  >{\raggedright\arraybackslash}p{(\linewidth - 4\tabcolsep) * \real{0.3333}}@{}}
\toprule\noalign{}
\begin{minipage}[b]{\linewidth}\raggedright
约分
\end{minipage} & \begin{minipage}[b]{\linewidth}\raggedright
完整名称（常见写法）
\end{minipage} & \begin{minipage}[b]{\linewidth}\raggedright
2100 分及以上：作用
\end{minipage} \\
\midrule\noalign{}
\endhead
\bottomrule\noalign{}
\endlastfoot
2300--2500 & 复杂数位动态规划（Advanced Digit Dynamic Programming） & 处理上界、前导零、数字和、禁止模式等多个数位状态。 \\
2200--2400 & 数位动态规划与自动机（Digit Dynamic Programming with Automaton） & 用自动机状态记录模式匹配、相邻限制和数位约束。 \\
1900--2200 & 带自环方程的概率动态规划（Probability Dynamic Programming with Self-Loops） & 把自环期望方程移项，或建立线性方程组后求解。 \\
2200--2500 & 概率动态规划与高斯消元（Probability Dynamic Programming with Gaussian Elimination） & 处理多个状态相互转移的期望和概率方程。 \\
1800--2200 & 子集枚举与所有子集和动态规划（Subset Enumeration and Sum Over Subsets Dynamic Programming，SOS DP） & --- \\
2100--2400 & 轮廓线动态规划（Profile Dynamic Programming） & 用一列/几列的状态描述窄网格的铺砖、放置和相邻约束。 \\
1800--2100 & 单调队列优化动态规划（Monotone Queue Optimization） & --- \\
2100--2400 & 斜率优化/凸包技巧（Convex Hull Trick，CHT） & 把二次或交叉项转成直线最值，利用斜率和查询单调性降复杂度。 \\
2200--2400 & 李超线段树优化动态规划（Li Chao Segment Tree Optimization） & 在线插入斜率不单调的转移直线并查询最优值。 \\
2200--2400 & 分治优化动态规划（Divide and Conquer Optimization） & 在决策点单调时缩小每层转移枚举范围。 \\
2300--2500 & 四边形不等式优化（Quadrangle Inequality Optimization） & 证明区间动态规划的决策单调性，从而减少转移枚举。 \\
2300--2500 & 决策单调性证明（Decision Monotonicity） & 证明最优决策点不向左移动，为分治优化或单调优化提供依据。 \\
2400--2500 & 动态动态规划与矩阵线段树 & 用矩阵乘积维护局部转移，支持单点修改后的全局答案。 \\
2400--2500 & 生成函数优化计数动态规划 & 把动态规划转成多项式运算，利用卷积或生成函数提速。 \\
2300--2500 & Steiner 树状压动态规划（Steiner Tree Subset Dynamic Programming） & 在关键点数量很少时求连接所有关键点的最小代价。 \\
2100--2400 & 树形背包优化（Tree Knapsack Optimization） & 合并子树选择数量，处理树上资源分配和选点限制。 \\
2300--2500 & 长链剖分优化树形动态规划（Long-Chain Decomposition for Tree DP） & 利用长链合并同深度状态，优化树形 DP 的空间或时间。 \\
\end{longtable}

\section{七、按题面快速查找}\label{ux4e03ux6309ux9898ux9762ux5febux901fux67e5ux627e}

\begin{itemize}
\tightlist
\item
  ``区间加/区间乘/区间和''：多重懒标记线段树或仿射标记线段树。
\item
  ``离线区间统计''：莫队算法；有修改时用带修改莫队算法；只能撤销时用回滚莫队算法。
\item
  ``三维偏序/时间和坐标双重限制''：三维偏序 CDQ 分治。
\item
  ``动态维护直线最值''：李超线段树；直线只在一段横坐标有效时用线段插入版李超线段树。
\item
  ``区间取最小/最大并查询区间和''：区间最值势能线段树。
\item
  ``树上子树颜色频次''：树上启发式合并；``动态树路径''：Link-Cut Tree 动态树。
\item
  ``模数不互质的多个同余式''：扩展中国剩余定理。
\item
  ``\passthrough{\lstinline!n!} 很大但递推阶数很小''：矩阵快速幂与线性递推。
\item
  ``所有倍数统计/恰好最大公约数''：莫比乌斯反演。
\item
  ``\passthrough{\lstinline!a\^x=b mod p!}''：Baby-step Giant-step 离散对数算法；底数与模数不互质时用扩展版本。
\item
  ``多项式卷积''：数论变换；``对称染色''：Burnside 引理或 Pólya 计数定理。
\item
  ``\passthrough{\lstinline!x²=a mod p!}''：Legendre 符号先判定，Tonelli--Shanks 或 Cipolla 算法求根。
\item
  ``不同子串/出现次数''：后缀自动机；``回文后缀链''：回文自动机；``多个模式串''：Aho--Corasick 多模式匹配自动机。
\item
  ``凸包最远点''：旋转卡壳；``多个线性不等式公共区域''：半平面交；``矩形并面积''：扫描线与线段树。
\item
  ``胜负态''：必胜态/必败态动态规划；``无偏且可拆分''：Sprague--Grundy 函数；``双方完整决策''：极大极小搜索与 Alpha-Beta 剪枝。
\end{itemize}

\section{八、常见算法组合}\label{ux516bux5e38ux89c1ux7b97ux6cd5ux7ec4ux5408}

\begin{enumerate}
\def\labelenumi{\arabic{enumi}.}
\tightlist
\item
  离散化 + 树状数组：排名、逆序、区间计数。
\item
  DFS 序 + 线段树：树上子树修改/查询。
\item
  最近公共祖先 + 树上差分：大量路径经过次数。
\item
  强连通分量缩点 + 有向无环图动态规划：有向图上的可达/最大权。
\item
  二分答案 + 广度优先搜索/贪心/网络流：判定型最优化。
\item
  Kruskal 重构树 + 倍增：按边权阈值查询连通块信息。
\item
  Aho--Corasick 自动机 + 动态规划：避免模式或统计匹配串。
\item
  可持久化线段树 + 最近公共祖先：树上路径第 \passthrough{\lstinline!k!} 小。
\item
  扫描线 + 线段树 + 坐标压缩：矩形并、覆盖次数。
\item
  回滚并查集 + 线段树时间分治：离线动态连通性。
\end{enumerate}

